\section{The Clock You Forgot You Were Using}
\label{sec:clocks-sysclk}

Every HAL has a constant called something like \codeinline{DEFAULT\_SYSCLK\_HZ}, and it is almost
always wrong for the board it runs on. This section is about the one on this board, because it
silently defines the length of every delay in the project.

\subsection{19 MHz, Not 78}

The EFR32MG24 can run at 78 MHz. It does not, out of reset, and this board's firmware never raises
it. The default is \codeinline{HFRCO0} at 19 MHz, and the measurement that settled it was a stopwatch
rather than a datasheet:

\begin{minted}{text}
# a blink with a 500 ms half-period, counted against a phone stopwatch
# 20 blinks took 20.0 s  ->  1.000 s per cycle  ->  the delay was right at 19 MHz
# the same code timed as if SYSCLK were 78 MHz produced a 4x fast blink
\end{minted}

That is a crude instrument, and it is the right one: a stopwatch over twenty cycles has an error far
below the 4× discrepancy being tested.

\subsection{Why the Error Is Invisible}

A wrong \codeinline{SYSCLK\_HZ} does not produce a wrong-looking program. It produces a program whose
delays are all scaled by the same factor — blinks that are too fast, tones that are too high, buses
that are too slow or too fast. Each of those looks like a bug in that peripheral:

\begin{table}[htbp]
    \centering
    \footnotesize
    \begin{tabular}{@{}p{0.30\textwidth}p{0.30\textwidth}p{0.32\textwidth}@{}}
        \toprule
        \textbf{Symptom} & \textbf{Blamed on} & \textbf{Actually} \\
        \midrule
        LED blinks at 4× the rate & the delay loop & \codeinline{SYSCLK\_HZ} too high \\
        Tone is 4× the intended pitch & the timer's TOP & same constant \\
        I2C bus works on one board, not another & the devices & same constant, other clock \\
        Buzzer sounds, then the beep pattern is wrong & the pattern code & same constant \\
        \bottomrule
    \end{tabular}
    \caption{One wrong constant, four different suspects. This is why it is measured once and written
    down next to the definition.}
    \label{tab:sysclk-symptoms}
\end{table}

\subsection{The Right Shape for the Constant}

The version that survives contact with a second board is not a number in the HAL. It is a number the
board passes in:

\begin{minted}{rust}
// The HAL knows how to delay; it does not know how fast this board's clock runs.
let mut delay = Delay::with_sysclk(syst, 19_000_000);
\end{minted}

and in this project's board layer, one place says it out loud:

\begin{minted}{rust}
pub const DEFAULT_SYSCLK_HZ: u32 = 19_000_000;   // measured, see 05.02
\end{minted}

A second board with a 38.4 MHz crystal changes one line, and every delay, tone and bus timing
follows it. A HAL with the constant baked in cannot be reused, and — worse — fails in a way that
looks like the new board's fault.

\subsection{Measuring It Properly}

The stopwatch is enough to catch a factor of four. For anything finer, use the cycle counter, which
counts CPU cycles and therefore cannot disagree with the CPU:

\begin{minted}{rust}
let start = dwt::cycles();
delay.delay_ms(1000);
let elapsed = dwt::since(start);
// elapsed / 19_000_000 must be ~1.0 s; a factor of 3.0 here means the delay's
// internal arithmetic is wrong, not the constant.
\end{minted}

That distinction — constant wrong \emph{versus} arithmetic wrong — is the difference between
recalibrating and debugging, and the cycle counter is what separates them.
